\chapter[Band 22: Gerichtete Graphen]{Band 22 auf einen Blick}
\noindent\textbf{Gerichtete Graphen}\par\medskip
Vorausgesetzt sind Relationen und endliche Folgen. Es gilt
\(\NatSeg n=\{i\in\mathbb N\mid i<n\}\); \(n\) zählt Kanten,
eine Knotenfolge hat daher \(n+1\) Einträge. Schleifen sind zugelassen.

\begin{BandUebersicht}
\textbf{Graph und Kantenrelation}
\UebFormel{\begin{aligned}
&\GraphStruct{V,E}\ \coloneqq\ E\subseteq V\times V.
\end{aligned}}
\UebQuellen{\FormulaRefAuto{GraphDef}[def]}
&
\textbf{Endpunkte liegen im Träger}
\UebFormel{\begin{aligned}
&\GraphStruct{V,E},\ (x,y)\in E\ \vdash\\
&\qquad x\in V\land y\in V.
\end{aligned}}
\UebQuellen{\FormulaRefAuto{GraphEdgeEndpointTyping}[thm]} \\
\addlinespace[.8ex]

\textbf{Ein- und Ausnachbarschaften}
Für \(\GraphStruct{V,E}\) und \(x\in V\):\UebFormel{\begin{aligned}
&\OutNeighbors E x=\{y\in V\mid(x,y)\in E\},\\
&\InNeighbors E x=\{y\in V\mid(y,x)\in E\}.
\end{aligned}}
\UebQuellen{\FormulaRefAuto{DirectedOutNeighborhoodDef}[def];
\FormulaRefAuto{DirectedInNeighborhoodDef}[def]}
&
\textbf{Mitgliedschaftskriterien}
\UebFormel{\begin{aligned}
&\GraphStruct{V,E},\ x,y\in V\ \vdash\\
&y\in\OutNeighbors E x\ \leftrightarrow\ (x,y)\in E,\\
&y\in\InNeighbors E x\ \leftrightarrow\ (y,x)\in E.
\end{aligned}}
\UebQuellen{\FormulaRefAuto{DirectedNeighborhoodMembership}[thm]} \\
\addlinespace[.8ex]

\textbf{Umkehrung und Einschränkung}
Für \(\GraphStruct{V,E}\) und \(U\subseteq V\):\UebFormel{\begin{aligned}
&\ReverseEdges E=\{(y,x)\mid(x,y)\in E\},\\
&\DirectedInducedEdges E U=E\cap(U\times U).
\end{aligned}}
\UebQuellen{\FormulaRefAuto{DirectedReverseEdgeRelationDef}[def];
\FormulaRefAuto{DirectedInducedEdgeRelationDef}[def]}
&
\textbf{Beide Konstruktionen erhalten Graphen}
\UebFormel{\begin{aligned}
&\GraphStruct{V,E}\ \vdash\
  \GraphStruct{V,\ReverseEdges E},\\
&\GraphStruct{V,E},\ U\subseteq V\ \vdash\\
&\qquad\GraphStruct{U,\DirectedInducedEdges E U}.
\end{aligned}}
\UebQuellen{\FormulaRefAuto{DirectedReverseGraph}[thm];
\FormulaRefAuto{DirectedInducedSubgraph}[thm]} \\
\addlinespace[.8ex]

\textbf{Walk und injektiver Pfad}
Bei festem \(\GraphStruct{V,E}\) bestehen die Walkaxiome aus\UebFormel{\begin{aligned}
&n\in\mathbb N,\quad p:\NatSeg{n+1}\to V,\\
&\forall i\in\NatSeg n\quad(p(i),p(i+1))\in E.\\
&\DirectedPath V E p n\ \leftrightarrow\\
&\quad\DirectedWalk V E p n
 \land p:\NatSeg{n+1}\inj V.
\end{aligned}}
\UebQuellen{\FormulaRefAuto{DirectedWalkDef}[def];
\FormulaRefAuto{DirectedPathDef}[def]}
&
\textbf{Typisierung und einzelne Schritte}
\UebFormel{\begin{aligned}
&\GraphStruct{V,E},\
 \DirectedWalk V E p n\ \vdash\\
&\qquad p(0)\in V\land p(n)\in V,\\
&\GraphStruct{V,E},\
 \DirectedWalk V E p n,\ i\in\NatSeg n\ \vdash\\
&\qquad(p(i),p(i+1))\in E.
\end{aligned}}
\UebQuellen{\FormulaRefAuto{DirectedWalkEndpointTyping}[thm];
\FormulaRefAuto{DirectedWalkStepEdge}[thm]} \\
\addlinespace[.8ex]

\textbf{Kreis und Kreisfreiheit}
Ein Kreis ist ein Walk positiver Länge mit\UebFormel{\begin{aligned}
&p(0)=p(n),\\
&\forall i,j\in\NatSeg n\
   (p(i)=p(j)\to i=j).
\end{aligned}}Kreisfreiheit schließt sämtliche solchen Kreise aus.
\UebQuellen{\FormulaRefAuto{DirectedCycleDef}[def];
\FormulaRefAuto{DirectedAcyclicGraphDef}[def]}
&
\textbf{Charakterisierung der Kreisfreiheit}
\UebFormel{\begin{aligned}
&\AcyclicGraphStruct{V,E}\ \eqvdash\\
&\quad\GraphStruct{V,E}\ \land\\
&\quad\neg\exists n\in\mathbb N\,\exists p\
          \DirectedCycle V E p n.
\end{aligned}}
\UebQuellen{\FormulaRefAuto{DirectedAcyclicGraphCharacterization}[thm]} \\
\addlinespace[.8ex]

\textbf{Erreichbarkeit}
Für \(u,v\in V\) bedeutet \(u\GraphReach E v\):\UebFormel{\begin{aligned}
&u=v\ \lor\ \exists n\in\mathbb N\,\exists p\\
&\quad(\DirectedWalk V E p n\\
&\qquad{}\land p(0)=u\land p(n)=v).
\end{aligned}}
\UebQuellen{\FormulaRefAuto{DirectedReachabilityDef}[def]}
&
\textbf{Reflexivität und Endpunkttypisierung}
\UebFormel{\begin{aligned}
&\GraphStruct{V,E},\ u\in V\ \vdash\
       u\GraphReach E u,\\
&\GraphStruct{V,E},\ u\GraphReach E v\ \vdash\\
&\qquad u\in V\land v\in V.
\end{aligned}}
\UebQuellen{\FormulaRefAuto{DirectedReachabilityReflexive}[thm];
\FormulaRefAuto{DirectedReachabilityEndpointTyping}[thm]} \\
\addlinespace[.8ex]
\end{BandUebersicht}

